\begin{afterword}
\summaryhead

The post quotes a New Scientist article. In it the cosmologist Max Tegmark defends the multiverse as a consequence of testable theories, and the philosopher Colin Howson proposes replacing Popper's falsificationism with a Bayesian account of science based on degrees of belief. The author says this shows the author to be less of a lonely iconoclast than it seems, and lists three differences from the mainstream effort to recast science in Bayesian terms.

First, being outside academia, the author can say ``extreme'' things that others may already think. Second, and chiefly, teaching probability theory will not be enough. It must be combined with research on cognitive biases and social psychology into a new art, ``Bayescraft''; the Bayesian theorist Robert Aumann, an Orthodox Jew, shows that probability theory alone does not do it. The author has not seen this suggested elsewhere. Third, it is possible to do much better than science: a Bayesian superintelligence could reach Newton's laws, and think general relativity worth testing, from the first falling apple it saw, while human minds, and academia even more, process evidence inefficiently.

\responsehead

The post corrects an impression the sequence had given. The previous week's posts set ``Science'' against ``Bayes'' without mentioning that a Bayesian philosophy of science already existed. Howson and Urbach's book on it first appeared in 1989, and Earman's in 1992. The post points to that literature and places the author within it.

The list of differences leaves out one that the quoted passage names. The article describes Howson's view as resting on ``the subjective concept of belief,'' and says Bayesian methods ``say nothing about what those initial beliefs should be.'' Howson is a subjective Bayesian: the Stanford Encyclopedia cites Howson for the view that Bayesian epistemology need not say what priors we should have, except that they be coherent. Two days earlier, ``Science Isn't Strict Enough'' denied the freedom that view allows: ``that doesn't mean you get a free, personal choice of making the priors whatever you want.'' On the question at the centre of that post, whether the evidence fixes one right estimate, Howson is on the other side. The agreement is real, but on this point narrower than the post's list suggests.

The main difference is offered as one the author has not ``seen suggested elsewhere.'' That is a claim about the author's reading. The idea had a named framework: by the 1980s, research on judgment and decision making paired normative models such as probability theory with descriptive models of how people err, and with prescriptive ``designs for improvement,'' which came to include teaching meant to counteract biases.

The example offered for this difference assumes that Aumann's religion is a failure of probabilistic reasoning. Aumann has described religion as ``an experience,'' ``mainly an emotional and aesthetic one,'' and as ``orthogonal'' to science.

The third difference, the superintelligence and the apple, is asserted here and argued later, in ``That Alien Message.''

In short: an acknowledgement that Bayesian philosophy of science is mainstream, whose list of differences leaves out a substantive one, the author's view that priors are not a free choice against Howson's subjective one, and which offers as unseen elsewhere a program the psychology of judgment had framed by the 1980s.
\end{afterword}
